% =====================================================================
%  Цель и задание лабораторной работы.
% =====================================================================

\textbf{Цель:} получить практические навыки установки, настройки и
применения системы обнаружения/предотвращения вторжений (IDS/IPS) для
защиты реального удалённого сервера, а также навыки сравнительного
анализа сетевой активности до и после её внедрения.

\vspace{1em}

\textbf{Задание}:
\begin{enumerate}
  \item Установить и настроить программное обеспечение, позволяющее
  усовершенствовать защиту сервера, согласно варианту.
  \item Выполнить анализ данных логов различных подсистем, в том числе
  SSH и веб-сервера удалённого сервера, за несколько дней.
  \item Оформить отчёт о проделанной работе.
\end{enumerate}

\textbf{Постановка задачи.} В качестве варианта выбран \texttt{OSSEC}
-- система обнаружения вторжений на основе хоста (HIDS) с
возможностью активного ответа (то есть способная работать и как IPS).
Она установлена на том же сервере (\texttt{test-lab4.work.gd}), что
использовался в лабораторной работе №4.1, и настроена на мониторинг тех
же подсистем -- SSH и веб-сервера \texttt{Caddy}, -- чтобы можно было
напрямую сравнить сетевую активность и её обработку до появления
OSSEC (данные лабораторной работы №4.1) и после включения его
активного ответа.
